\subsection{Separately concave residuals and minimum cuts}\label{sec:cuts}

In this subsection
\[
F(x)=c_0+\sum_i(a_ix_i^2+b_ix_i)+\sum_{i<j}c_{ij}x_ix_j
\]
is a rational quadratic. Its coordinates are split into a \emph{core}
$\mathcal C$ of $k$ continuous coordinates and a \emph{residual}
$\mathcal R$. Missing couplings have $c_{ij}=0$. Assume
\begin{equation}\label{eq:cr-class}
\text{(R1)}\ a_i\le0\ (i\in\mathcal R),\qquad
\text{(R2)}\ c_{ij}o_io_j\le0\ (i,j\in\mathcal R)\ \text{for some }
o\in\{-1,1\}^{\mathcal R}.
\end{equation}
Nothing is assumed about core--core or core--residual coefficients, and
the residual interaction graph may be dense. A residual coordinate may be
continuous or integer; integer bounds are rounded inward first and an
empty domain is rejected. Given the signs of the residual $c_{ij}$, a
breadth-first traversal finds an orientation satisfying (R2) or a cycle
proving that none exists.

\begin{lemma}[Endpoint reduction]\label{lem:cr-endpoint}
Fix $v\in\bar X_{\mathcal C}$ and a box $\prod_{i\in\mathcal R}[l_i,u_i]$
(after inward rounding). Under (R1), the minimum of $F(v,\cdot)$ over the
continuous box is attained at a point with $x_i\in\{l_i,u_i\}$ for all
$i\in\mathcal R$. Consequently it equals the minimum over the mixed
residual domain.
\end{lemma}
\begin{proof}
Start from a minimizer, which exists by compactness. For each residual
coordinate in turn, the objective as a function of that coordinate alone
is $a_it^2+(\text{affine in }t)$, which is concave by (R1); replace the
coordinate by an endpoint with no larger value. The result is an endpoint
minimizer. Its coordinates are integral for integer coordinates, so it is
feasible for the mixed domain, which is contained in the continuous box.
\end{proof}

For $i\in\mathcal R$ put $\alpha_i=l_i$, $d_i=u_i-l_i$ if $o_i=1$, and
$\alpha_i=u_i$, $d_i=-(u_i-l_i)$ if $o_i=-1$. For $y\in\{0,1\}^{\mathcal R}$
let $x_i(y)=\alpha_i+d_iy_i$; this enumerates all residual endpoint
vectors. In the network below, for a node set $S$ that contains the
source $\mathsf s$ but not the sink $\mathsf t$, $\operatorname{cap}(S)$
denotes the total capacity of arcs leaving $S$.

\begin{proposition}[Signed cut representation]\label{prop:cr-cut}
Fix a rational core point $v$. Under~\eqref{eq:cr-class} there are rationals
$\varphi_0,\mu_i,\omega_{ij}$, computable by exact evaluation of $F$, such
that for all $y\in\{0,1\}^{\mathcal R}$
\[
F(v,x(y))=\varphi_0+\sum_{i}\mu_iy_i+\sum_{i<j}\omega_{ij}y_iy_j,
\qquad \omega_{ij}=c_{ij}d_id_j\le0 .
\]
Let $\rho_i=\mu_i+\frac12\sum_{j\ne i}\omega_{ij}$ and
$r_{ij}=-\omega_{ij}/2\ge0$. Build a network on
$\mathcal R\cup\{\mathsf s,\mathsf t\}$ with arcs $i\to\mathsf t$ of
capacity $\max\{\rho_i,0\}$, arcs $\mathsf s\to i$ of capacity
$\max\{-\rho_i,0\}$, and, for each pair with $\omega_{ij}\ne0$, arcs
$i\to j$ and $j\to i$ of capacity $r_{ij}$. With
$S_y=\{\mathsf s\}\cup\{i:y_i=1\}$,
\begin{equation}\label{eq:cr-cutid}
F(v,x(y))=\varphi_0+\sum_i\min\{0,\rho_i\}+\operatorname{cap}(S_y).
\end{equation}
\end{proposition}
\begin{proof}
Expand $F(v,\alpha+d\circ y)$. Core terms and $c_0$ give constants. A term
$a_ix_i^2$ gives $a_i\alpha_i^2+(2a_i\alpha_id_i+a_id_i^2)y_i$ by
$y_i^2=y_i$; $b_ix_i$ and core--residual terms $c_{ij}v_jx_i$ are affine
in $y_i$; and $c_{ij}x_ix_j$ for residual $i,j$ gives
$c_{ij}d_id_jy_iy_j$ plus affine terms. By (R2),
$\omega_{ij}=c_{ij}o_io_j(u_i-l_i)(u_j-l_j)\le0$. For binary $y$,
$\omega y_iy_j=\frac\omega2(y_i+y_j)-\frac\omega2|y_i-y_j|$, so
\[
F(v,x(y))=\varphi_0+\sum_i\rho_iy_i+\sum_{i<j}r_{ij}|y_i-y_j| .
\]
In the cut $S_y$, an arc $i\to\mathsf t$ is cut exactly when $y_i=1$, an
arc $\mathsf s\to i$ exactly when $y_i=0$, and exactly one of the two pair
arcs exactly when $y_i\ne y_j$. Hence
$\operatorname{cap}(S_y)=\sum_i[\max\{\rho_i,0\}y_i+\max\{-\rho_i,0\}(1-y_i)]
+\sum_{i<j}r_{ij}|y_i-y_j|$, and
$\rho_iy_i=\max\{\rho_i,0\}y_i+\max\{-\rho_i,0\}(1-y_i)+\min\{0,\rho_i\}$.
\end{proof}

\begin{theorem}[Certified cut recourse]\label{thm:cr-oracle}
Under~\eqref{eq:cr-class}, for every rational core point $v$, every
rational residual subbox and every rational change of the residual linear
coefficients, one maximum-flow computation returns an exact conditional
minimizer, the exact value $W_{\mathcal C}(v)$, and a certificate: an
endpoint label $y$ and a feasible flow whose value equals
$\operatorname{cap}(S_y)$. The bit complexity is polynomial in the
encoding length of the query, with an absolute exponent independent of
$k$. Given the network, the certificate is checked with
$O(|\mathcal R|+m_{\mathcal R})$ rational operations, where $m_{\mathcal R}$
is the number of nonzero residual couplings. Conditions (R1)--(R2) are
preserved by residual subboxes, by fixing residual coordinates, by added
linear terms and by fixing the core.
\end{theorem}
\begin{proof}
By Lemma~\ref{lem:cr-endpoint} the conditional minimum is attained at an
endpoint label, and by~\eqref{eq:cr-cutid} minimizing over labels is a
minimum $\mathsf s$--$\mathsf t$ cut problem with nonnegative capacities.
Any feasible flow has value at most the capacity of any cut (weak
duality), so a flow whose value equals $\operatorname{cap}(S_y)$ proves
that $S_y$ is a minimum cut; checking it needs capacity bounds,
conservation at each residual node and one comparison. A maximum flow and
such a cut exist and are found, for instance, by the Edmonds--Karp
algorithm with $O(|V||A|^2)$ arithmetic operations, where
$|V|=|\mathcal R|+2$ and $|A|=O(|\mathcal R|+m_{\mathcal R})$. After
multiplying all capacities by a common denominator (whose bit length is
polynomial in the query), capacities are integers, every augmentation is
by an integer, and every flow value is an integer bounded by the total
capacity. Thus all intermediate numbers have polynomial bit length. The
coefficients $\varphi_0,\mu_i,\omega_{ij}$ are obtained by exact
rational evaluation. The stability statement holds because (R1)--(R2)
involve only $a_i$ and $c_{ij}$ for residual $i,j$, which these operations
do not change; a fixed coordinate has $d_i=0$.
\end{proof}

\begin{remark}[Scope and prior work]\label{rem:cr-cut-prior}
Each ingredient is classical: endpoint optimality of coordinatewise
concave functions; the representation of quadratic pseudo-Boolean
functions with nonpositive couplings by cuts
\cite{Ivanescu1965,PicardRatliff1975,KolmogorovZabih2004}; and the resulting
polynomial solvability of box QPs that are submodular with nonpositive
diagonal, as recorded in \cite{BurerNatarajanWillemsen2026v3}. The role of
Theorem~\ref{thm:cr-oracle} here is a box-stable, certified conditional
oracle with an arbitrary core. Separate concavity (R1) is essential:
continuous submodular box QP without it is not known to be polynomial,
and its natural SDP relaxation is exact only up to dimension three
\cite{BurerNatarajanWillemsen2026v3}. Del Pia and Khajavirad also separate
nonpositive-diagonal variables, which can be taken binary, from the others
and prove polynomial classes with few positive-diagonal variables; their
conditions bound the treewidth and interfaces of the binary part rather
than its signs \cite{DelPiaKhajavirad2026}.
\end{remark}

\paragraph{Search over a continuous core.}\label{sec:cr-search}

Let the core domain be $[0,1]^k$ ($k\ge1$, continuous) and the residual
domain $X_{\mathcal R}$ any compact mixed box. Put
$V(v)=W_{\mathcal C}(v)=\min_{z\in X_{\mathcal R}}F(v,z)$ and let $L\ge0$ be an
upper coordinate curvature in the core coordinates (for a quadratic,
$L\ge\max_{i\in\mathcal C}2a_i$). Assume an \emph{exact oracle}: for every
dyadic $v\in[0,1]^k$ it returns $V(v)$ and a completion $z(v)\in X_{\mathcal R}$
with $F(v,z(v))=V(v)$, with a certificate. Theorem~\ref{thm:cr-oracle}
supplies one under~\eqref{eq:cr-class}.

\paragraph{Core search.}
Level $j$ uses dyadic cells of side $h_j=2^{-j}$ and
$e_j=kLh_j^2/8$. Start with $\mathcal Q_0=\{[0,1]^k\}$ and $U=+\infty$.
At level $j$: query $V$ at all corners of the cells of $\mathcal Q_j$ that
were not queried before; set $U$ to the least value of all completions
found so far; let
$\lambda_j=\min_{C\in\mathcal Q_j}(\min_{\text{corners of }C}V-e_j)$;
retain $\mathcal K_j=\{C\in\mathcal Q_j:\min_{\text{corners}}V-e_j\le U\}$;
and let $\mathcal Q_{j+1}$ be the $2^k$ dyadic children of each retained
cell. Stop when $U-\min\{\lambda_j,U\}\le\varepsilon$ or at a level limit.

\begin{theorem}[Core search]\label{thm:cr-search}
For every level $j$:
\begin{enumerate}
\item[(a)] every cell of $\mathcal Q_j$ containing $x^*_{\mathcal C}$ for a
global optimizer $x^*$ is retained; in particular $\mathcal K_j\ne\emptyset$;
\item[(b)] $\lambda_j\le\OPT\le U\le\lambda_j+e_j$;
\item[(c)] every retained cell has a corner $v$ with $V(v)\le\OPT+2e_j$;
\item[(d)] every $x\in X$ satisfies $F(x)\ge\min\{\lambda_j,U\}$, and this
can be verified from the list of cells, the certified queried values and
the retention decisions, without any optimization.
\end{enumerate}
\end{theorem}
\begin{proof}
Core coordinates are continuous, so every cell of $\mathcal Q_j$ has
$e_{\mathcal C}(C)=e_j$. (a) The unit cube contains $x^*_{\mathcal C}$. If
$C\in\mathcal Q_j$ contains it, Lemma~\ref{lem:cr-cell} gives
$\min_{\text{corners}}V-e_j\le\OPT\le U$, so $C$ is retained and the child
containing $x^*_{\mathcal C}$ belongs to $\mathcal Q_{j+1}$.
(b)--(c) This is Theorem~\ref{thm:cr-filter} with $\ell=V$ and $\eta=0$;
its hypothesis on an optimizer cell holds by (a).
(d) Children of retained cells cover them, so the cells of $\mathcal Q_j$
together with the cells discarded at levels $i<j$ cover $[0,1]^k$. If
$x_{\mathcal C}$ lies in a cell of $\mathcal Q_j$, then
$F(x)\ge V(x_{\mathcal C})\ge\lambda_j$ by Lemma~\ref{lem:cr-cell}. If it lies in
a cell discarded at level $i$, then
$F(x)\ge\min_{\text{corners}}V-e_i>U^{(i)}\ge U$, where $U^{(i)}$ is the
incumbent at level $i$. Each step uses only recorded values.
\end{proof}

For any $\varepsilon>0$ the search therefore stops by the first level with
$e_j\le\varepsilon$. Its cost is controlled by the number of near-optimal
grid points. Let $N_j$ be the number of points
$v\in h_j\Z^k\cap[0,1]^k$ with $V(v)\le\OPT+2e_j$.

\begin{lemma}[Charging queries]\label{lem:cr-charge}
Level $0$ makes $2^k$ queries and level $j+1$ makes at most $8^kN_j$.
\end{lemma}
\begin{proof}
By Theorem~\ref{thm:cr-search}(c), every retained level-$j$ cell has a
corner counted in $N_j$, and a grid point is a corner of at most $2^k$
cells; so $|\mathcal K_j|\le2^kN_j$. Each retained cell has $2^k$
children with $2^k$ corners each.
\end{proof}

\begin{proposition}[Deterministic count under core growth]\label{prop:cr-growth}
Suppose $V(v)-\OPT\ge g\norm{v-v^*}^2$ for all $v\in[0,1]^k$, some $v^*$
and some $g>0$; this holds with the same $g$ if $F$ has point growth $g$.
Let $\kappa=\max\{1,L/g\}$. Then levels $0,\ldots,J$ make at most
\[
2^k+8^kJ\bigl(1+\sqrt{k\kappa}\bigr)^k
\]
queries. The constant $g$ is not used by the algorithm.
\end{proposition}
\begin{proof}
If $F(x)-\OPT\ge g\norm{x-x^*}^2$, then minimizing over the residual gives
$V(v)-\OPT\ge g\norm{v-x^*_{\mathcal C}}^2$. A point counted in $N_j$
satisfies $g\norm{v-v^*}^2\le2e_j=kLh_j^2/4$, so each coordinate lies
within $\frac{h_j}2\sqrt{k\kappa}$ of $v^*_i$. An interval of length
$h_j\sqrt{k\kappa}$ contains at most $1+\sqrt{k\kappa}$ points of
$h_j\Z$. Hence $N_j\le(1+\sqrt{k\kappa})^k$; apply
Lemma~\ref{lem:cr-charge} and sum over $j<J$.
\end{proof}

\begin{example}[Flat directions]\label{ex:cr-flat}
Without growth the count can be exponential in the number of accuracy
bits. Take $k=2r$, no residual, and $V(v)=\sum_{t=1}^r(v_{2t-1}-v_{2t})^2$,
so $L=2$ and $\OPT=0$. Every product of diagonal cells
$\prod_t\bigl([a_th_j,(a_t+1)h_j]\times[a_th_j,(a_t+1)h_j]\bigr)$ has a
corner with $V=0$, hence is retained. So level $j$ retains at least
$2^{jr}=2^{jk/2}$ cells, and gap $\varepsilon$ costs
$\Omega((L/\varepsilon)^{k/4})$ queries.
\end{example}

Random linear perturbations of the core costs remove this effect in
expectation: if the core coefficients are drawn independently and uniformly
from $M\ge2^{J-1}$ equally spaced points of $[-\sigma,\sigma]$, the search is
correct for every draw and the expected number of queries through level $J$
is at most $2^k+8^kJ[3+(1+k/2)L/(2\sigma)]^k$ (Theorem~\ref{thm:cr-smoothed}
in Appendix~\ref{app:smoothed}). The bound is linear in the number of levels,
whereas Example~\ref{ex:cr-flat} is exponential in it, and the residual
problem enters only the cost of each query.

\paragraph{Exact output.}\label{sec:cr-exact}

\begin{lemma}[Uniform height over residual labels]\label{lem:cr-height}
Let $F$ be a rational quadratic with continuous core $[0,1]^k$ and residual
satisfying (R1), with integer residual bounds rounded inward. Let
$\Lambda_c$ be the least common denominator of all coefficients of $F$,
$\Lambda_e$ that of all residual bounds, $T_0=\Lambda_c\Lambda_e^2$,
$H_0=\max\{1,\max_{i,j\in\mathcal C}|T_0\,\partial_{ij}F|\}$ and
\[
D_0=T_0\bigl(k!\,H_0^k\bigr)^2 .
\]
For every residual endpoint label $z$, the value
$\Phi(z)=\min_{v\in[0,1]^k}F(v,z)$ is a rational with reduced denominator
at most $D_0$, and so is $\OPT$.
\end{lemma}
\begin{proof}
Fix $z$ and put $\phi(v)=F(v,z)$. Every monomial of $T_0\phi$ has an
integer coefficient: the coefficients of $v_iv_j$ and $v_i^2$ are
coefficients of $F$; the coefficient of $v_i$ is $b_i+\sum_jc_{ij}z_j$;
and the constant collects $c_0$, $a_jz_j^2$, $b_jz_j$ and $c_{ij}z_iz_j$
for residual $i,j$. Each has denominator dividing $\Lambda_c\Lambda_e^2$.
The matrix $A=T_0\nabla^2_{\mathcal C\mathcal C}F$ is integral with entries
bounded by $H_0$ in absolute value.

Among minimizers of $\phi$ on $[0,1]^k$ choose one, $v$, with the fewest
coordinates in $(0,1)$; let $S$ be that set. The other coordinates are
$0$ or $1$. If $S=\emptyset$, then $T_0\phi(v)\in\Z$. Otherwise
$\nabla_S\phi(v)=0$ and $A_{SS}\succeq0$ by first- and second-order
optimality on the face. If $A_{SS}$ were singular, a nonzero $w$
supported on $S$ with $A_{SS}w_S=0$ would give $\phi(v+tw)=\phi(v)$ for all
$t$, and moving $t$ until a coordinate of $S$ reaches $0$ or $1$ would
produce a minimizer with fewer interior coordinates. Hence
$\delta=|\det A_{SS}|$ is a positive integer, and by the Leibniz expansion
$\delta\le|S|!\,H_0^{|S|}\le k!\,H_0^k$. The stationarity system
$A_{SS}v_S=-T_0\nabla_S\phi(0_S,v_{\mathcal C\setminus S})$ has an integral
right-hand side, so $\delta v$ is integral by Cramer's rule. Evaluating
the integer-coefficient quadratic $T_0\phi$ at $v=(\delta v)/\delta$ gives
an element of $\delta^{-2}\Z$. Thus $\Phi(z)\in(T_0\delta^2)^{-1}\Z$.
By Lemma~\ref{lem:cr-endpoint}, some global optimizer has an endpoint
label $z^*$, and $\OPT=\Phi(z^*)$.
\end{proof}

\begin{remark}
Any common multiple $S_0$ of all denominators can replace $\Lambda_c$ and
$\Lambda_e$, for example $T_0=S_0^3$. A cube, or the mixed form
$\Lambda_c\Lambda_e^2$, is needed in general: the term $\frac12x_ix_j$ at
residual endpoints $x_i=x_j=\frac12$ equals $\frac18$, while the least
common denominator of the data is $2$.
\end{remark}

\begin{theorem}[Exact output with a cut residual]\label{thm:cr-exact}
Assume~\eqref{eq:cr-class} with continuous core $[0,1]^k$ and rational data.
Run the core search until $U-\min\{\lambda_j,U\}<D_0^{-2}$, which happens
by the first level with $e_j<D_0^{-2}$. Let $z_U$ be the residual label
of the incumbent completion. Minimize $\phi(v)=F(v,z_U)$ exactly over
$[0,1]^k$ by enumerating the $3^k$ faces of the cube, solving the
stationarity system on each face whose free Hessian block is nonsingular,
and keeping feasible solutions. The best candidate $\hat v$ satisfies
$F(\hat v,z_U)=\OPT$, so $(\hat v,z_U)$ is a global optimizer. A checker
needs only the search trace, $D_0$, feasibility and exact value of the
output, its reduced denominator ($\le D_0$), and the inequality
$F(\hat v,z_U)-\min\{\lambda_j,U\}<D_0^{-2}$.
If in addition $V$ has core growth with constant $g$ as in
Proposition~\ref{prop:cr-growth}, the total bit complexity is
$f(k,\kappa)(I+1)^{C_1}$ with an absolute constant $C_1$, and $g$ is not
an input.
\end{theorem}
\begin{proof}
The proof of Lemma~\ref{lem:cr-height} shows that a minimizer of $\phi$
lies on some face with a nonsingular free block, so the enumeration
contains it; all candidates are feasible, so the best one attains
$\Phi(z_U)$. If $(v_U,z_U)$ is the incumbent completion, then
$\OPT\le\Phi(z_U)\le F(v_U,z_U)=U$ and, by Theorem~\ref{thm:cr-search}(d),
$\min\{\lambda_j,U\}\le\OPT$. So $0\le\Phi(z_U)-\OPT<D_0^{-2}$. Both
numbers have reduced denominators at most $D_0$ by
Lemma~\ref{lem:cr-height}, and distinct such rationals differ by at least
$D_0^{-2}$; hence they are equal. The checker's test proves optimality of
any feasible point with these properties by the same separation argument.

For the complexity, $\log T_0$ and $\log H_0$ are $O(I)$ and $k\le I$, so
$\log D_0=O(I^2)$ and the number of levels is
$J=O(\log D_0+\log(kL)+1)=\poly(I)$. Under core growth,
Proposition~\ref{prop:cr-growth} bounds the queries by
$2^k+8^kJ(1+\sqrt{k\kappa})^k$. Each query is a cut problem whose data have
bit length polynomial in $I+J$, hence $\poly(I)$ by
Theorem~\ref{thm:cr-oracle}. Face enumeration costs $3^k\poly(I)$. All
factors depending only on $k$ and $\kappa$ form $f$.
\end{proof}

Without core growth the exact mode still terminates, but its number of
cells can be $2^{kJ}$ with $J=\Theta(\log D_0)$. The stopping depth must
not be combined with Theorem~\ref{thm:cr-smoothed} by choosing
$M\ge2^{J}$: sampled coefficients enter $D_0$, and for a fixed noise scale
the atom denominators grow with $M$, so this condition generally cannot
hold.

